\documentclass[pdflatex,sn-basic,Numbered]{sn-jnl}% Basic Springer Nature Reference Style, numbered citations (Wireless Personal Communications house style)
\usepackage{amsmath,amssymb}
\usepackage{graphicx}
\usepackage{booktabs}
\usepackage{xcolor}
\usepackage{algorithm}
\usepackage{algorithmic}

% Visible TODO marker for content the author (not Claude) still needs to
% fill in or decide -- same convention as Papers_4-5/Paper_4/main.tex.
\newcommand{\authorTODO}[1]{\par\noindent\fbox{\parbox{0.95\linewidth}{\textbf{\color{red}TODO(MANOJ):} #1}}\par}

% M49b-0: visible draft flags marking content under live re-collection
% (docs/PAPER5_M48b_provenance_audit.csv found the entire Live Single-gNB
% Anchor section ran pre-M41-fix, below the scheduler floor on all three
% slices). Two forms: a boxed flag for section-level content, an inline
% one for a single flagged sentence embedded in flowing prose (abstract,
% Introduction) where a full box would be disruptive. Neither deletes the
% flagged text -- it is the re-collection target, not discarded.
\newcommand{\prefixflag}[1]{\par\noindent\fbox{\parbox{0.95\linewidth}{\textbf{\color{red}PRE-FIX DATA -- DO NOT SUBMIT (M49b re-collection):} #1}}\par}
\newcommand{\prefixinline}[1]{\textcolor{red}{[PRE-FIX DATA -- M49b re-collection: #1]}}

\begin{document}

\title[GAT-CTDE for Coordinated Multi-gNB Admission Control]{GAT-CTDE: A Graph-Attention Architecture for Coordinated Multi-gNB Slice Admission Control}

\author[1]{\fnm{Manoj Prasad} \sur{Kunasegran}}\nomail

\author*[1]{\fnm{Wai Leong} \sur{Pang}}\email{TODO@taylors.edu.my}

\author[1]{\fnm{Swee King} \sur{Phang}}\nomail

\affil*[1]{\orgdiv{School of Engineering}, \orgname{Taylor's University}, \orgaddress{\city{Subang Jaya}, \country{Malaysia}}}

\abstract{Coordinating slice admission decisions across multiple gNBs
is a genuinely multi-agent problem: does a shared, graph-structured
cluster representation help per-gNB agents make better joint decisions
than a single agent or independent learners? We study this with
GAT-CTDE, a centralised-training, decentralised-execution architecture
built on a shared Graph Attention encoder feeding decentralised
per-slice Q-heads, extended to federated training with differential
privacy and evaluated under mid-episode disruption. Scaled from 3 to
19 gNBs, GAT-CTDE resists a training collapse that becomes total for a
flattened-state single-agent baseline, converging, across 64
independently-drawn seeds, on a residual 35.4\% [25.5\%, 45.8\%]
collapse rate. Federating the architecture costs nothing measurable in
reward, and differential privacy degrades decision precision only past
a real noise threshold. Under gNB dropout, demand spikes, and agent
churn, its disruption cost accelerates with severity like every arm
tested, with a coordination-specific penalty only at the mildest
severity. On the metric that reflects the actual training objective,
GAT-CTDE decisively beats an independent-learner ablation, though its edge over the single-agent baseline remains small and not statistically decisive at this sample size. This is an honest limit we report rather than obscure. These results were made trustworthy, not just favourable, by
two safeguards this paper also contributes: a correctness-aware metric
pair that caught a real training collapse the standard SLA-compliance
metric scored no worse than correct behaviour, and a
reproduction-versus-replication discipline that caught and retracted
one architectural overclaim under disruption before publication. A
live single-gNB anchor confirms the architecture's decision logic
transfers to real traffic unchanged. \prefixinline{this live anchor ran before the M41 scheduler-floor fix; re-collection pending, see docs/PAPER5\_M48b\_provenance\_audit.csv}}

\keywords{Network Slicing, Slice Admission Control, Multi-Agent Reinforcement Learning, Graph Attention Networks, Federated Learning, Open RAN}

\maketitle

\noindent\footnotesize
M.~P.~Kunasegran: orcid.org/0009-0005-3169-1873.
W.~L.~Pang: orcid.org/0000-0001-8407-5648.
S.~K.~Phang: orcid.org/0000-0002-7877-8766.
\normalsize
\vspace{6pt}

\section{Introduction}
Network slicing lets a single physical radio access network serve heterogeneous service classes, namely Ultra-Reliable Low-Latency Communication (URLLC), enhanced Mobile Broadband (eMBB), and massive Machine-Type Communication (mMTC), under distinct per-slice Service Level Agreements (SLAs). Our earlier work~\cite{mecon-paper1,icraie-paper2}
established that a Deep Q-Network (DQN) admission controller can learn
to enforce these SLAs under a single gNodeB (gNB), validated on a live
OpenAirInterface (OAI) testbed. A companion submission~\cite{cacs26-paper} extends this single-gNB live validation to a QoE-aware reward variant, independently confirming the live E2 control loop's feasibility (median round-trip latency 0.57\,ms) for a policy under 20,000 parameters. A real multi-cell deployment does not admit requests at one gNB in isolation, since neighbouring cells share backhaul, spectrum planning, and a cluster-wide load-balance objective, making coordinated admission a genuinely multi-agent problem. This raises a direct question: does a shared, graph-structured representation of the cluster's state help a set of per-gNB agents make better joint decisions than a single agent using the flattened joint state, or independent per-gNB agents with no shared representation at all?

This paper's contribution is GAT-CTDE, a centralised-training,
decentralised-execution multi-agent architecture built on a shared
Graph Attention Network~\cite{gat} encoder over the gNB topology
feeding decentralised per-slice Q-heads, and the evidence for how it
scales, coordinates, and degrades under realistic operating
conditions. Extended to federated training with differential privacy and evaluated under mid-episode gNB dropout, demand spikes, and agent churn, we scale the architecture to cluster sizes $N \in \{7, 19\}$ under ring and hex-grid adjacency. It resists a training collapse that becomes total for a flattened-state single-agent baseline, converging on a residual 35.4\% collapse rate across 64 independently-drawn seeds rather than either of two disagreeing small-sample estimates. On the metric that reflects the actual training objective, the architecture decisively beats an independent-learner ablation that isolates the same coordination contribution, though its edge over the single-agent baseline specifically is small and not statistically decisive at this sample size. This is an honest limit we report rather than smooth over to fit a cleaner architecture story. Federating the architecture costs nothing
measurable in reward, and differential privacy degrades decision
precision only past a real threshold rather than a smooth trade-off;
under disruption, its cost accelerates with severity like every arm
tested, with a coordination-specific penalty at only the mildest
severity tested.

Reporting these results honestly required two methodological
safeguards we also contribute, since neither is standard practice in
this literature. First, the standard episode-level SLA-compliance metric can score a degenerate, always-accept policy no worse than genuinely differentiated decision-making, because blocking a low-priority slice under congestion is the reward-optimal action, yet it costs compliance without ever recovering it. We introduce a correctness-aware metric pair (mean reward per step and block precision, both computed from data every run already logs) that exposes this gap, and use it to diagnose and fix a real training collapse our first encoder suffered on all 30 campaign seeds. Second, we establish a reproduction-versus-replication discipline. Every headline result is checked against an independent seed sample rather than only reproduced on the same seeds, a distinction that caught and retracted a real architectural overclaim under mid-episode disruption before it could be reported as a finding. We finally anchor the architecture's
single-gNB baseline against a real OAI testbed, confirming its
offline-trained decision logic transfers to live traffic unchanged. \prefixinline{pre-M41-fix live anchor, re-collection pending}

\section{Related Work}
\label{sec:litreview}
Our own systematic review~\cite{ref:access} screened 6{,}286 records down to 67 core 2024--2026 studies on machine-learning-driven, user-centric distributed network slicing, organised around three enabling pillars, namely topology-aware graph representation, distributed multi-agent control, and privacy-preserving federation, and found individual pillars well developed but rarely combined under realistic constraints, and never with a validated live deployment. None of the works we reviewed, there or here, checks whether its own evaluation metric can mask a degenerate policy, or replicates a headline finding on an independent seed sample. This is the gap this paper fills.

The canonical result behind our own encoder choice is
Shao~\emph{et al.}~\cite{shao2021gatmarl}, who show a graph attention
encoder strengthens inter-agent cooperation over plain DQN/A2C
baselines in dense-cellular slicing. Closest to our admission-control setting, Ahmadi~\emph{et al.}~\cite{ahmadi2024generalizable} generalise coordinated multi-agent slicing across unseen topologies via graph attention, a topology-generalisation question distinct from, but complementary to, the evaluation-integrity question we ask on a fixed topology. We test the complementary direction directly in Section~\ref{sec:results-m6}, asking whether correctness-aware metrics remain reliable as cluster size and adjacency sparsity change, rather than whether a policy generalises to an unseen topology. Sulaiman~\emph{et
al.}~\cite{sulaiman2023coordinated} is the closest prior art to our own
problem, multi-agent DRL jointly solving slicing and admission control
across the substrate network; neither that work nor its
topology-generalisation follow-up reports whether its own
reward-optimal policy could be mistaken for a degenerate one, or
checks a headline result against an independent seed sample. Tashman and Cherkaoui~\cite{tashman2026adversarial} and Fatehi~\emph{et al.}~\cite{fatehi2026attentionmarl} both study robustness to a disruptive condition, as our disruption campaign (Section~\ref{sec:results-m4}) does; neither checks whether the metric reporting that robustness is itself trustworthy, or replicates a robustness finding on an independent sample. This is exactly the gap our own disruption campaign's independent-seed replication closes, and not hypothetically, since it is what overturned our own initial churn-immunity reading.

On the federated side, Zhang~\emph{et
al.}~\cite{zhang2022federatedoran} coordinate independent O-RAN xApps
via federated DRL, the closest precedent to our own periodic-
synchronisation motivation, but without a differential-privacy step or
a correctness-aware metric to check whether federating changed
decision quality. Yasin~\emph{et al.}~\cite{yasin2025dpfl} combine federated edge learning with differential privacy to secure the O-RAN RIC (the closest FL+DP precedent) but target a different threat model (securing the controller) rather than characterising a control policy's own privacy-utility trade-off, which our $\sigma$-sweep (Section~\ref{sec:results-m3}) reports. A small further set of works
begins to combine two of these three
pillars~\cite{mohajer2026pronto,bouroudi2025gnnmarl,qasim2025tgnnfl,gong2026femaddpg},
but, like every work surveyed here, evaluates exclusively on its own
reported outcome metric, computed once, on one seed or run. This is the gap this paper fills, not a new point in the GAT/MARL/FL design space but a demonstration (on GAT-CTDE, extended to federation, privacy, and disruption) that checking a metric's own trustworthiness and replicating a finding on an independent sample are both necessary, and that skipping either changes what a reader is entitled to believe the resulting numbers mean.

\section{System Model and Problem Formulation}
\label{sec:system-model}
We consider a cluster of $N{=}3$ gNBs, each serving the same three
slices (URLLC, eMBB, mMTC) under a shared per-gNB Physical Resource
Block (PRB) budget $B$. An admission decision is realised as a
PRB-ceiling adjustment (\texttt{slicing\_control\_m} min/max ratio),
not a literal per-flow accept/reject, since the real OAI xApp control
API only exposes ratio ceilings. No real multi-gNB physical topology is
available to us, so we take the least assumption-laden choice and
treat the cluster as fully connected, noting this explicitly as a
modelling choice rather than a measured fact. At each control step,
every gNB has zero or more pending admission requests, each tagged
with a slice identity, and an agent decides accept/reject for each
independently, conditioned on the current cluster state; the discrete
action space matches our single-gNB work exactly
(\texttt{action\_dim=2}), so any coordination benefit found here is
attributable to the architecture, not a richer action space.

The reward combines a per-slice SLA-compliance term with a
cluster-wide load-balance penalty, weighted by a per-slice
\texttt{priority\_weight} (URLLC highest, mMTC lowest) and a congestion
coefficient tuned, in prior work, so that the reward-optimal policy
under congestion is \emph{differentiated shedding}: continue accepting
URLLC, mostly accept eMBB, and reject mMTC. This is confirmed directly by Q-value inspection on a working baseline policy (Section~\ref{sec:collapse}), not assumed. Every result in this paper
except Section~\ref{sec:results-m8}'s bounded live anchor is obtained
in an offline, closed-loop simulation environment deliberately
oversubscribed to roughly 110\% of one gNB's PRB budget, reusing our
existing multi-gNB configuration unmodified. A separate investigation
found this offline environment's demand process does not reproduce the
live testbed's SLA-margin distribution closely enough to support a
live-rank prediction claim, even after recalibrating the demand
model's temporal structure; we therefore use it explicitly as a \emph{stress environment} for a contention regime our live rig does not reach, not as a claim that any offline ranking would reproduce live. Section~\ref{sec:results-m8} tests this directly rather than only arguing for it. Multi-gNB claims in this paper remain
offline-only throughout; this rig cannot yet sustain a live multi-gNB
test to check them against.

\section{GAT-CTDE Method}
\label{sec:methodology}

\begin{figure}[t]
\centering
\includegraphics[width=0.92\textwidth]{figures/fig1_gat_ctde_fl_architecture.pdf}
\caption{GAT-CTDE architecture. (a) A shared Graph Attention encoder
over the fully-connected 3-gNB topology produces one embedding per
node; per-slice Q-heads, each with its own parameters, consume that
node's own embedding row at action-selection time (decentralised
execution). (b) The federated training variant: each gNB is a client
with a private copy of the same network, synchronised via periodic
FedAvg with an optional per-client DP-SGD clip-and-noise step.}
\label{fig:architecture}
\end{figure}

\subsection{Shared Encoder, Decentralised Q-Heads}
Fig.~\ref{fig:architecture}(a) shows the architecture. At every control
step, each gNB's own per-slice [PRB-used-ratio, congestion-level,
queue-length-norm] features form that node's row in a joint
node-feature matrix. A two-layer, four-head GAT~\cite{gat}, implemented
directly (the small graph sizes here make a dense, masked-attention
implementation as fast as a sparse one), produces one embedding per
gNB. Centralisation comes from this encoder, since its parameters are updated by every agent's loss jointly, so it sees the full cluster's state during training. Decentralised execution comes from what happens next. At action-selection time, each agent consumes only its own node's embedding row to decide its own pending requests.

\label{sec:methodology-perslice}
An early version fed each embedding, concatenated with a one-hot slice
context, into a single shared Q-head MLP. Under this design, mMTC's
true reward-optimal margin under congestion is small by the reward's
own calibration, while URLLC's much larger priority weight produces
proportionally larger gradients; because all three slices update the \emph{same} shared parameters, URLLC's abundant, large-magnitude gradients overwrite mMTC's much smaller true signal before it registers. This is a standard shared-parameter/class-imbalance failure mode.
We replace the single shared head with \emph{one small MLP per slice}, selected by the context one-hot where the head is applied (Fig.~\ref{fig:architecture}(a)). Each slice's decision boundary now has its own parameters, so one slice's gradients can no longer directly overwrite another's.

Algorithm~\ref{alg:gat-ctde} states the training loop, in which the encoder is updated from every request in a sampled minibatch regardless of gNB or slice (centralised training), while action selection at rollout time consults only the acting agent's own embedding row (decentralised execution). Every pending request in a step is its own independent
temporal-difference sample, never summed or averaged across a step
before the loss is computed, matching our existing single-agent DQN's
own granularity.

\begin{algorithm}[t]
\caption{GAT-CTDE: centralised training, decentralised execution}
\label{alg:gat-ctde}
\begin{algorithmic}[1]
\REQUIRE shared GAT encoder $f_\theta$, per-slice Q-heads $\{Q_{\phi_k}\}_{k=1}^{K}$, target networks $f_{\theta^-}, \{Q_{\phi_k^-}\}$, replay buffer $\mathcal{D}$, fully-connected adjacency $A$, discount $\gamma$, target-sync period $C$ episodes
\FOR{episode $= 1, 2, \dots$}
  \STATE reset environment; observe joint node features $X_1 \in \mathbb{R}^{N \times d}$
  \FOR{step $t = 1, \dots, T$}
    \STATE $E_t \gets f_\theta(X_t, A)$ \COMMENT{one embedding row per gNB; centralised: sees every node}
    \FOR{each pending request $r$ at gNB $i$, slice $k$}
      \STATE $a_r \gets \arg\max_a Q_{\phi_k}(E_t[i])[a]$ w.p.\ $1{-}\epsilon$, else uniform random \COMMENT{decentralised: own row $E_t[i]$ only}
    \ENDFOR
    \STATE $X_{t+1}, \rho_t, \text{done} \gets \textsc{env.step}(\{a_r\})$
    \STATE store $\big(X_t, X_{t+1}, \{(i,k,a_r)\}_r, \rho_t, \text{done}\big)$ in $\mathcal{D}$
    \IF{$|\mathcal{D}| \geq$ warmup size}
      \STATE sample minibatch $B \subset \mathcal{D}$
      \STATE $E_B \gets f_\theta(X_B, A)$; \ $E'_B \gets f_{\theta^-}(X'_B, A)$ \COMMENT{target net, no grad}
      \FOR{each request $(i,k,a) \in B$}
        \STATE $q \gets Q_{\phi_k}(E_B[i])[a]$
        \STATE $y \gets \rho + \gamma \cdot (1{-}\text{done}) \cdot \max_{a'} Q_{\phi_k^-}(E'_B[i])[a']$
      \ENDFOR
      \STATE $\mathcal{L} \gets$ mean Huber loss over every request's $(q, y)$ pair \COMMENT{per-request, never summed per step}
      \STATE update $\theta$, $\{\phi_k\}$ by $\nabla \mathcal{L}$; clip grad-norm to 1.0
    \ENDIF
  \ENDFOR
  \IF{episode $\bmod\ C = 0$}
    \STATE $\theta^- \gets \theta$; \ $\phi_k^- \gets \phi_k\ \forall k$
  \ENDIF
\ENDFOR
\end{algorithmic}
\end{algorithm}

\subsection{Federated Training with Differential Privacy}
Fig.~\ref{fig:architecture}(b) shows the federated variant, in which each gNB becomes an FL client with a private copy of the same network, training only on its own requests. Transitions are partitioned by owning agent, so no client sees another client's raw admission data, while its GAT encoder still consumes the full joint node-feature snapshot every arm observes. Every $R$ local steps, the server averages all
clients' online weights (FedAvg~\cite{fedavg}) and broadcasts the
result back. We also implement FedProx~\cite{fedprox} (a proximal term pulling local weights toward the last broadcast global model), motivated by an implicit, seeded per-gNB load multiplier this project's own cluster-scaling investigation (Section~\ref{sec:results-m6}) later revealed as genuine, uncontrolled client heterogeneity. Swept at $\mu \in \{0.01, 0.1, 1.0\}$ under both a homogeneous-load control and this heterogeneity, FedProx earns no measurable dividend at any tested strength. This traces to a concrete mechanism, since real per-round client drift within \texttt{local\_steps\_per\_round} averages only 0.03\% of the TD loss, too small for the proximal term to matter. A follow-up at 10$\times$ the round length lets real drift grow nearly 50-fold (to a mean 1.44\%) yet still produces eval-time behaviour bit-identical to FedAvg, ruling out round length as the limiting factor rather than leaving it untested.

Differential privacy is applied as a standard DP-SGD~\cite{dpsgd} step, in which every client's gradient is clipped to a fixed norm and, when enabled, perturbed with Gaussian noise scaled by a noise multiplier $\sigma$, before that client's own optimiser step. We report $\sigma$ as the
primary privacy parameter, alongside a formal $(\varepsilon,\delta)$
bound computed via zero-concentrated differential privacy
(zCDP)~\cite{zcdp} composition of the Gaussian mechanism ($\rho =
k/(2\sigma^2)$ over $k$ per-client releases,
$\varepsilon(\delta)=\rho+2\sqrt{\rho\ln(1/\delta)}$). This is a conservative, non-subsampled bound, reported alongside $\sigma$ rather than in place of it. Algorithm~\ref{alg:fed-dp} states one federated round precisely;
clipping is applied identically whether or not noise is added, so a $\sigma{=}0$ run and a $\sigma{>}0$ run differ only in whether noise is injected. This is the isolation the privacy-cost comparison (Section~\ref{sec:results-m3}) depends on.

\begin{algorithm}[t]
\caption{Federated GAT-CTDE round with DP-SGD}
\label{alg:fed-dp}
\begin{algorithmic}[1]
\REQUIRE $N$ clients, each with local network $(\theta_i, \{\phi_{i,k}\})$; global snapshot $\bar\theta$; local steps per round $R$; clip norm $C_{\text{dp}}$; noise multiplier $\sigma$; FedProx weight $\mu$ ($=0$ unless noted)
\FOR{round $= 1, 2, \dots$}
  \FOR{local step $= 1, \dots, R$, each client $i$ independently}
    \STATE sample minibatch $B_i$ of $i$'s own requests from $i$'s local buffer
    \STATE $E \gets f_{\theta_i}(X_{B_i}, A)$ \COMMENT{full joint state; client's OWN local weights only}
    \STATE $\mathcal{L}_i \gets$ per-request Huber TD loss (Alg.~\ref{alg:gat-ctde}, lines 11--13) $+\ \frac{\mu}{2}\lVert \theta_i - \bar\theta \rVert^2$ \COMMENT{FedProx term}
    \STATE $g_i \gets \nabla_{\theta_i,\phi_i} \mathcal{L}_i$
    \IF{$\sigma > 0$}
      \STATE $g_i \gets \text{clip}(g_i,\, C_{\text{dp}})$; \ $g_i \gets g_i + \mathcal{N}(0,\, (\sigma C_{\text{dp}})^2 I)$ \COMMENT{DP-SGD~\cite{dpsgd}}
    \ELSE
      \STATE $g_i \gets \text{clip}(g_i,\, C_{\text{dp}})$ \COMMENT{same clip norm, no noise -- isolates the privacy cost}
    \ENDIF
    \STATE $(\theta_i, \phi_i) \gets (\theta_i, \phi_i) - \eta\, g_i$
  \ENDFOR
  \STATE $\bar\theta, \bar\phi \gets \frac{1}{N}\sum_{i=1}^N (\theta_i, \phi_i)$ \COMMENT{FedAvg~\cite{fedavg}, equal client weights}
  \FOR{each client $i = 1, \dots, N$}
    \STATE $(\theta_i, \phi_i) \gets (\bar\theta, \bar\phi)$ \COMMENT{broadcast}
  \ENDFOR
\ENDFOR
\end{algorithmic}
\end{algorithm}

\subsection{Experimental Setup}
Both the centralised and federated campaigns use 300 training and 50
evaluation episodes per seed. The centralised campaign uses a fixed
30-seed list (seeds 900--929) with three arms: GAT-CTDE; an
independent-DQN ablation (one per-gNB DQN instance, no parameter
sharing, no topology information) that isolates the GAT/centralised-
training contribution; and a single-agent DQN baseline run against the
flattened joint state. All three arms use matched hyperparameters
(discount factor $\gamma=0.95$, per-episode epsilon decay). The
federated campaign uses the first 10 of these 30 seeds and a
five-level noise-multiplier sweep $\sigma \in \{0, 0.5, 1, 2, 4\}$, so
its $\sigma{=}0$ level is directly comparable to the centralised
campaign's same 10 seeds. All confidence intervals are 95\% bootstrap
percentile intervals (10{,}000 resamples); paired comparisons use the
Wilcoxon signed-rank test.

\section{Evaluation-Integrity Findings: A Training Collapse, Diagnosed and Fixed}
\label{sec:collapse}
This section diagnoses a training collapse that a standard
SLA-compliance metric could not detect on its own. We rely throughout
on two correctness-aware metrics computed from data every run already
logs: mean reward per step, the actual training objective, and block
precision, the fraction of blocked requests that correctly target the
slice the reward calibration makes reject-optimal.
Section~\ref{sec:correctness-metrics} below states why these are
necessary in full; here they let us tell a genuine training failure
apart from a legitimate one.

While generating a privacy-utility sweep for the federated variant
(Section~\ref{sec:results-m3}), noise levels that should have produced
genuinely different, stochastically-trained checkpoints instead
produced bit-identical per-seed evaluation compliance. Tracing this back, we found that our first GAT-CTDE encoder's eval policy blocked zero requests across every evaluation episode, on all 30 campaign seeds. An always-accept policy is a functional no-op, so any two checkpoints that reach it produce identical trajectories regardless of their internal weights. We confirmed, rather than assumed, that this was a genuine failure and not a legitimate optimum. A direct Q-value probe on the collapsed checkpoint showed $Q(\text{accept}) - Q(\text{reject})$ positive and large for every slice, including mMTC, while the independent-DQN ablation (Section~\ref{sec:results-m2}), which sees identical raw features through no encoder at all, does not exhibit this collapse. This isolates the encoder, not the reward or training loop, as the fault.

We compared the encoder's own embedding for a synthetic idle state against a congested one on the collapsed checkpoint. The raw input difference has $\ell_2$ norm 2.55 and the resulting embedding difference has norm 19.6, an apparently large, congestion-sensitive response, but the cosine similarity between the two embeddings is 0.99997. This means the encoder was encoding congestion almost entirely as embedding \emph{magnitude}, with no mechanism preventing this and no rotation of direction. Since the downstream Q-head is roughly linear in its input, larger magnitude pushes $Q(\text{accept})$ up faster than $Q(\text{reject})$ as congestion increases, which runs backwards from the reward's intent.

We add a per-layer \texttt{LayerNorm} to the encoder, constraining every node's embedding to zero-mean, unit-variance regardless of input magnitude. This is a real, verified improvement (collapse drops from 30/30 to 27/30 seeds) but not a complete fix. We then apply the per-slice
Q-heads of Section~\ref{sec:methodology-perslice}, which reduces
collapse further, to 9/30 seeds (Fig.~\ref{fig:collapse}); a further variant disabling LayerNorm's learnable affine transform made results \emph{worse} (0/3 pilot seeds differentiated versus 1/3 with it enabled) and was reverted. Not every plausible fix helped, and we report the one that did not.

\begin{figure}[t]
\centering
\includegraphics[width=\columnwidth]{figures/fig2_collapse_reduction.pdf}
\caption{Fraction of the 30 campaign seeds reaching genuine,
correctly-targeted differentiated shedding (mMTC-only blocking under
congestion), across the original unnormalised encoder and the two
fixes described in Section~\ref{sec:collapse}.}
\label{fig:collapse}
\end{figure}

\label{sec:eval-log-bug}
A separate, unrelated bug was caught while building an evaluation-time disruption harness. Our logging utility opens its output file in append mode and never truncates it, and GAT-CTDE was retrained in place three times while diagnosing the collapse above, each time writing evaluation results to the same per-seed path. Every one of the 30 seeds' evaluation logs had silently accumulated all three runs' episodes stacked together, confirmed directly by each log's own episode index resetting to 1 twice. Every GAT-CTDE number reported
through an earlier draft of this paper, and every M3 result reusing
GAT-CTDE's centralised evaluation as a reference point, was therefore
a hidden three-way average across two architecturally-obsolete runs
and the true final one. The saved model checkpoints themselves were
unaffected (a different save routine overwrites rather than appends,
verified by reproducing one seed's final-block statistics exactly from
its checkpoint in isolation), so the fix required no retraining: we
cleared every stale log, re-ran evaluation once against the
already-trained checkpoints, and hardened the training script to
clear its output path before any future run. All numbers in this paper
reflect that corrected data.

\label{sec:correctness-metrics}
The fixes above make the architecture more correct while making its \texttt{sla\_compliance\_all\_slices} score \emph{worse} (Section~\ref{sec:results-m2}), because blocking mMTC under congestion, the reward-optimal action, does not recover mMTC's own SLA margin in this stress regime, so every newly-correct decision costs compliance the metric never credits back. We therefore report two further metrics,
both computed from data every arm already logs: \textbf{mean reward
per step}, the actual training objective; and \textbf{block
precision}, the fraction of blocked requests that target mMTC
specifically, the only slice the reward calibration ever makes
reject-optimal (undefined, and reported as such, for seeds that never
block anything). Any admission-control evaluation where correct behaviour cannot recover an already-stressed slice's compliance should report a decision-level correctness metric alongside, not instead of, an outcome-level compliance one. This is the general lesson this specific collapse makes concrete.

\section{Multi-gNB Results: Centralised, Federated, and Disruption}

\begin{figure}[t]
\centering
\includegraphics[width=\textwidth]{figures/fig3_m2_results.pdf}
\caption{M2 centralised campaign, per-seed evidence for the
correctness-aware metrics (Section~\ref{sec:correctness-metrics}). (a)
Paired per-seed comparison against single-agent DQN on mean reward per
step. Each line is one seed, coloured by whether GAT-CTDE scored
higher (green, 20), tied (grey, 3), or lower (red, 7), the source of
this figure's borderline $p=0.0577$. (b) Full per-seed reward
distribution for all three arms; GAT-CTDE is split into seeds that
learned genuine differentiated shedding (blue, 21) versus seeds still
stuck at the always-accept collapse (grey $\times$, 9), shown without
claiming the split predicts reward, since several still-collapsed
seeds score among the arm's highest.}
\label{fig:m2-results}
\end{figure}

\subsection{Centralised Campaign}
\label{sec:results-m2}

The $N=3$ fully-connected topology used throughout this campaign is
deliberately the setting where a shared graph-attention encoder has
the least topology to exploit, since adjacency across three gNBs is
close to trivial at full connectivity. This makes it a clean substrate
for the evaluation-integrity findings that follow, uncomplicated by
any topology effect; the encoder's coordination advantage is put to a
genuine test only once the cluster scales to $N=19$ under sparser
adjacency in Section~\ref{sec:results-m6}.

Table~\ref{tab:m2-results} and Fig.~\ref{fig:m2-results} report the
final, post-fix 30-seed campaign. On \texttt{sla\_compliance\_all\_slices},
GAT-CTDE's mean is 0.151, below its own pre-fix number (0.353) since
the metric penalises newly-correct mMTC-blocking decisions, and the paired comparison against single-agent DQN shows no clear direction, at +0.036 (95\% CI [$-$0.050, 0.126], Wilcoxon $p=0.4549$, 11 wins / 9 ties / 10 losses). On mean reward per step, the metric that reflects the actual objective, GAT-CTDE beats the independent-DQN ablation decisively, by +1.461 (95\% CI [0.550, 2.782], $p<0.0001$, 27/0/3), isolating a real contribution from the shared GAT encoder and centralised training, but its edge over single-agent DQN is small and not statistically decisive at this sample size, at +0.195 (95\% CI [$-$0.017, 0.442], $p=0.0577$, 20/3/7). Its block precision (0.952) is high regardless, since when GAT-CTDE blocks a request, it is almost always the reward-correct slice. 21 of the 30 seeds show genuine, correctly-targeted differentiated shedding, up from 3/30 under the LayerNorm-only fix and 0/30 originally. This is a real, substantial behavioural improvement that does not, at this sample size, translate into a decisive reward edge over the single-agent baseline specifically.

\begin{table}[t]
\caption{M2 campaign results, 30 seeds/arm, 95\% bootstrap CIs.}
\label{tab:m2-results}
\centering
\begin{tabular}{lccc}
\toprule
Arm & Mean reward/step & Block precision & Compliance$^{\ast}$ \\
\midrule
GAT-CTDE & \textbf{14.062} [13.91, 14.23] & 0.952 & 0.151 [0.083, 0.225] \\
Indep. DQN & 12.601 [11.30, 13.49] & 0.901 & 0.049 [0.013, 0.093] \\
Single-agent & 13.867 [13.69, 14.05] & 0.987 & 0.115 [0.061, 0.176] \\
\botrule
\end{tabular}
\begin{tablenotes}
\footnotesize
\item[$\ast$] Low compliance is the metric artifact Section~\ref{sec:collapse} explains, not a performance failure.
\end{tablenotes}
\end{table}

\begin{table}[t]
\caption{Training hyperparameters, all fixed across seeds.}
\label{tab:hyperparams}
\centering
\begin{tabular}{ll}
\toprule
Hyperparameter & Value \\
\midrule
Learning rate (GAT-CTDE) & $1{\times}10^{-4}$ \\
Learning rate (DQN baselines) & $1{\times}10^{-3}$ \\
Discount $\gamma$ & 0.95 \\
Epsilon: start $\to$ end (decay/episode) & $1.0 \to 0.05$ (0.985) \\
Target-network sync & every 10 episodes \\
Optimiser & Adam \\
Training / evaluation episodes & 300 / 50 \\
\midrule
\multicolumn{2}{l}{\emph{GAT encoder (GAT-CTDE only)}} \\
Hidden / output dim & 16 / 16 \\
Attention heads & 4 \\
Encoder layers & 2 \\
Q-head hidden dim & 32 \\
\midrule
\multicolumn{2}{l}{\emph{Federated + DP (federated GAT-CTDE only)}} \\
Local steps per round & 50 \\
Aggregator & FedAvg ($\mu$=0) \\
DP gradient clip norm & 1.0 \\
DP noise multiplier sweep $\sigma$ & $\{0, 0.5, 1, 2, 4\}$ \\
\botrule
\end{tabular}
\end{table}

\subsection{Federated Training and Differential Privacy}
\label{sec:results-m3}

\begin{figure}[t]
\centering
\includegraphics[width=\textwidth]{figures/fig4_m3_privacy.pdf}
\caption{M3 federated and differential-privacy sweep, per-seed
evidence. (a) Paired per-seed federation cost on mean reward per step,
same 10 seeds. Centralised scores higher on 2 seeds, federated on 3, 5
ties, no measurable difference (Wilcoxon $p=0.8125$). (b)
Privacy-utility curve using block precision against the DP noise
multiplier $\sigma$, individual seed values (light dots) behind the
mean, with the seed count that ever blocked anything at each $\sigma$
printed below (precision is undefined otherwise). Precision is perfect
through $\sigma=1.0$, dropping sharply by $\sigma=4.0$.}
\label{fig:m3-privacy}
\end{figure}

Comparing the centralised campaign against the federated variant with no DP noise, on the same 10 seeds, mean reward per step moves from 14.073 to 13.940, a change of +0.133 (95\% CI [$-$0.138, 0.470], Wilcoxon $p=0.8125$). This is no measurable cost to federating at this sample size, and both arms' block precision is perfect (1.000) on seeds that block anything at all. Sweeping the DP noise multiplier (Fig.~\ref{fig:m3-privacy}(b)), block precision is perfect through $\sigma=1.0$ then drops sharply, to 0.834 at $\sigma=2.0$ and 0.584 at $\sigma=4.0$. This is a real threshold, not a smooth degradation, while mean reward per step stays in a comparatively narrow band across the whole sweep (12.45--14.19), since DP noise erodes \emph{which} slice gets blocked, not whether the policy tries to differentiate at all. The compliance-based version of this same sweep (not shown) is non-monotonic with sign flips between adjacent $\sigma$ levels. This is a second instance, independent of Section~\ref{sec:collapse}'s architectural finding, of \texttt{sla\_compliance\_all\_slices} failing to track the quantity we actually care about.

\subsection{Disruption Resilience}
\label{sec:results-m4}
We evaluate every arm's already-trained, frozen checkpoint (no
retraining) against three mid-episode disruptions injected at
evaluation time: \emph{gNB dropout} (one gNB's requests are forced to
reject and its row in the joint feature matrix is zeroed, for a
window); \emph{demand spike} (a transient arrival-rate multiplier for a
fixed window); and \emph{agent churn} (one gNB's decisions are made by a fresh, never-trained copy of its own policy class instead of the loaded checkpoint, for a window). Agent churn is a different, well-motivated failure mode from dropout (active-but-uncoordinated vs.\ absent) that applies to every genuinely multi-agent arm; single-agent DQN is excluded.
Severity is the window's duration (10\%/30\%/60\% of the episode) for
dropout and churn, and the arrival-rate multiplier
(2$\times$/4$\times$/8$\times$) for spike. We normalise reward by each
step's own pending-request count before averaging, since raw reward
mechanically inflates with request volume regardless of decision
quality.

\begin{figure}[t]
\centering
\includegraphics[width=\textwidth]{figures/fig5_m4_disruption.pdf}
\caption{M4 disruption campaign, 10 seeds/arm. (a)-(b) Normalised reward cost (baseline $-$ disrupted) vs.\ severity for dropout and agent churn. Every arm, every severity, shows a real cost that accelerates with severity rather than growing linearly. (c) Block precision vs.\ demand-spike multiplier: flat and near-ceiling for every arm, unlike (a)-(b), since spike does not threshold on the metric that matters.}
\label{fig:m4-disruption}
\end{figure}

Each disruption kind intercepts the frozen policy's per-step
evaluation loop differently: dropout acts twice, corrupting the target
gNB's input row and separately forcing that gNB's own requests to
reject, so a policy cannot simply learn around the corrupted feature;
agent churn substitutes a fresh, never-trained policy's actions for
the target gNB only, leaving every other agent's decisions untouched;
demand spike is entirely external to the policy, acting on the
arrival process the environment synthesises for the following step.
All three are evaluation-only perturbations, and the policy itself is never updated.

Dropout costs every arm significantly at every severity (Wilcoxon $p\leq0.03$, 9 or 10 of 10 seeds hurt in every case), and the cost accelerates with severity rather than growing proportionally. For GAT-CTDE, the cost is +0.165 at the 10\% window, +0.601 at 30\%, and +1.710 at 60\%, and the same super-linear pattern holds for every arm. Agent churn shows
the same accelerating shape for GAT-CTDE (+0.059/+0.200/+0.529,
$p=0.0039$ throughout) and the federated arm (+0.073/+0.240/+0.585,
$p=0.0039$ throughout). Independent DQN's own churn cost needs both samples reported together to state honestly. On the seeds used throughout the rest of this paper (900--909), it is small and not conventionally significant (+0.000/+0.168/+0.640, $p=0.084/0.106/0.065$), which an earlier pass read as independent DQN being architecturally immune to churn, since it never attends to another agent's features. An independent-seed replication (seeds 1000--1009, disjoint from every other result in this paper) contradicts that reading directly. Independent DQN's churn cost there is significant at every severity (+0.097/+0.657/+1.223, $p=0.0020$ throughout) and \emph{larger} than GAT-CTDE's or the federated arm's at matched severities, the opposite of immune. We do not conclude independent DQN is immune to churn, since the direction was consistent across both samples (churn always hurts, never helps). What was wrong was treating a borderline, underpowered $n=10$ result as a confirmed architectural property, and this section reports the corrected reading, not the retracted one, as this paper's discipline of checking every headline result against an independent seed sample, not the same seeds reproduced again, requires.

Demand spike shows no threshold at all on block precision, which stays
flat and near-ceiling through the highest tested multiplier for every
arm; a reward-based signal that spike \emph{improves} per-step reward
is real and reproduces across both seed samples for GAT-CTDE,
independent DQN, and the federated arm ($p\leq0.014$ in both), but traces to a fixed per-step violation penalty diluted by a larger request-volume denominator, not better decision-making. Block precision, a volume-invariant ratio, is the metric that actually answers whether spike corrupts decision quality, and it says no.
Single-agent DQN's version of this same signal is significant in the original sample ($p=0.0020$) but not the replication ($p=0.084$); the direction is consistent, the significance is sample-dependent, and both are reported as such.

Every dropout and churn result above was independently checked against the 1000--1009 sample and holds up in direction and approximate magnitude. For example, GAT-CTDE's dropout costs there are +0.170/+0.552/+1.449, close to the 900--909 sample's +0.165/+0.601/+1.710. Every headline paired comparison in this paper
was checked this way, recomputed from raw logs on both samples:
GAT-CTDE's edge over independent DQN (Section~\ref{sec:results-m2})
holds decisively on the replication sample too (+0.767, 95\% CI
[0.391, 1.233], $p<0.0001$, versus the committed sample's +1.461
[0.550, 2.782]); its edge over single-agent DQN stays non-decisive on
both (+0.445 [$-$0.049, 1.280], $p=0.111$, versus +0.195 [$-$0.017,
0.442], $p=0.058$); federation's reward cost stays statistically
indistinguishable from zero on both, with the point estimate's sign
flipping between samples ($-$0.080 [$-$0.241, 0.001] on replication
versus +0.133 [$-$0.138, 0.470] committed) but neither ever
significant. Of the seven headline paired comparisons checked this way across the whole paper, six show overlapping confidence intervals across samples; the one exception, independent DQN's churn cost, is the retraction above.

A formal paired Wilcoxon test on (GAT-CTDE's cost $-$ independent DQN's cost), computed across the same 10 seeds for every (kind, severity) cell, reaches significance only once, at dropout at the mildest severity (10\% window), where GAT-CTDE's cost is significantly \emph{higher} than independent DQN's (+0.0135, 95\% CI [0.0053, 0.0236], Wilcoxon $p=0.0059$, 9 losses / 0 ties / 1 win). Every other cell, dropout at 30\%/60\% and churn at all three severities, shows no significant difference ($p \geq 0.23$). We report the result as it is, real at one severity and absent everywhere else tested, not a general claim that coordination costs more under disruption.

\section{Cluster-Size Scaling}
\label{sec:results-m6}
We extend the cluster from $N=3$ to $N \in \{7, 19\}$ under three
adjacency structures (fully-connected, ring, hex-grid), scaling
synthetic arrivals with $N$ to hold offered load per gNB roughly
constant. GAT-CTDE is the only arm whose Q-head consumes this
adjacency; independent DQN and single-agent DQN never read a graph
structure and serve as a topology-invariance check. We normalise
reward per gNB before averaging (the same fix Section~\ref{sec:results-m4}
applies on the severity axis) and report block precision alongside,
since \texttt{sla\_compliance\_all\_slices} is OR-aggregated across gNBs and trends toward 0 as $N$ grows regardless of policy quality. This is a third instance of this paper's central evaluation-integrity finding, now on the cluster-size axis. The primary sample is 12 seeds
(900--911) at $N=19$, all three topologies, all three arms; an
independent-seed replication (1000--1002, disjoint) checks
generalisation.

At the full 12-seed sample, the paired per-gNB reward-margin comparison at $N=19$ is not significant for any topology (fully-connected $-$0.175 [95\% CI $-$0.528, +0.076], $p=0.7422$; ring $-$0.704 [$-$1.879, +0.042], $p=0.7422$; hex $-$0.695 [$-$1.873, +0.053], $p=0.9102$). This is mechanically explicable, not a red flag, since single-agent DQN's always-accept collapse (below) collects the reward's service term on every request with no rejection cost, while GAT-CTDE pays a real cost every time it correctly rejects mMTC. This
paper's evidence for GAT-CTDE's advantage at $N=19$ lives in collapse
resistance and block precision, not the reward-margin metric.

\begin{figure}[t]
\centering
\includegraphics[width=\columnwidth]{figures/fig6_m6_topology.pdf}
\caption{$N=19$ cluster, per-seed evidence across independent seed samples. (a) Collapse rate (fraction of eval runs with zero blocks), pooled across all three topologies. Single-agent DQN collapses totally in every sample; GAT-CTDE's rate varies sample-to-sample (per-sample breakdown in the body text), converging to a combined, seed-level-bootstrapped estimate of 35.4\% [25.5\%, 45.8\%] across all 64 independent seeds (diamond marker, error bar). (b) Block precision by topology, primary sample, non-collapsed seeds only (single-agent DQN has none to show). GAT-CTDE's non-collapsed seeds are mostly near-ceiling but with a low-precision seed at every topology, while independent DQN sits in a narrow, mediocre band, identical across topologies as expected for an arm that never consumes adjacency.}
\label{fig:m6-topology}
\end{figure}

Single-agent DQN collapses totally at $N=19$: 12/12 primary-sample cells and 3/3 replication cells, every topology, zero exceptions (15/15 combined). This is confirmed at the individual-decision level by a direct Q-value probe identical in method to Section~\ref{sec:collapse}'s original diagnosis. GAT-CTDE collapses only partially, and independent
DQN never fully collapses (0/36) but never cleanly differentiates
either, sitting in a narrow, mediocre 0.29--0.49 block-precision band
identical across all three topologies, as expected for an arm that
never consumes adjacency (Section~\ref{sec:results-m27} revisits this specific claim under a recalibrated simulator and finds a small, rare exception to "never"). GAT-CTDE's own non-collapsed seeds show a topology-dependent secondary failure mode, since most hold near-perfect precision, but at least one seed per topology drops to 0.0--0.5, including one seed showing perfect precision (1.000) at fully-connected and near-zero (0.000) at ring/hex on the identical seed, environment stream, and initial weights.

GAT-CTDE's own collapse rate initially disagreed sharply between the
primary sample (31\%) and the replication sample (78\%); rather than
report the range as open, we drew a third, independent sample targeted
specifically at this question (49 seeds, disjoint from both prior
samples). Pooling all three samples gives 68/192 cells collapsed, 35.4\% (seed-level bootstrap 95\% CI [25.5\%, 45.8\%] across the 64 independent seeds, since collapse status correlates within a seed across its three topologies). GAT-CTDE's collapse rate at $N=19$ is best understood as roughly 35\% with substantial seed-to-seed variability, and none of the three samples was wrong, since each was an honest reading of a genuinely wide distribution, and no single small sample should be trusted alone.

\section{Live Single-gNB Anchor}
\label{sec:results-m8}
\prefixflag{Every live run in this section (both anchor runs below) used \texttt{saclb\_live.yaml} between commits 82b20f6 (2026-07-16) and 40ae2f3 (2026-08-19) -- before the M41 scheduler-floor fix (cda9d65, 2026-09-05). M42's config-floor audit confirms all three slices' ratio ranges sit below the scheduler's 5-PRB grant minimum for this exact config. The decision-logic-transfer claim below may be unaffected, revised, or overturned once re-collected post-fix; see docs/PAPER5\_M48b\_provenance\_audit.csv and docs/PAPER5\_M49b\_supervisor\_note.md.}
Prior live-testbed work on this rig already found that the offline
environment cannot predict live \emph{ranking} (Spearman
$\rho=0.097$, $p=0.855$, across six single-gNB checkpoints); we do not
repeat that attempt, and instead ask a narrower, bounded question this
rig can answer: does the offline-trained \emph{decision logic} survive
contact with real traffic at all. This rig cannot yet sustain a live
multi-gNB test, so GAT-CTDE's multi-gNB coordination benefit and the
federated/DP results are architecturally untestable on it; we trained a fresh single-agent-DQN checkpoint dimension-matched to the live evaluation path and ran it against a real OAI gNB, 3 UEs, and real per-slice UDP traffic, twice, once from a leftover backlog state and once from a verified-clean baseline.

Both live runs blocked only mMTC (46/32 and 45/32 requests, zero urllc or eMBB blocks in either), exactly matching this checkpoint's offline block precision of 1.000. The offline-trained differentiated-shedding decision transfers to live traffic essentially unchanged, replicated across two independent runs. Both runs' SLA-margin readings, by contrast, diverged from anything offline produces. eMBB's margin fell from $+1.0$ to $-1{,}002{,}377.5$ within a 10-minute run, recurring to the decimal in the independent second run. This traces to a concrete unit mismatch, not left unexplained, since the margin computation divides a per-slice backlog reading by a fixed calibration constant tuned for the offline environment's synthetic proxy, but live, the identically-named field is wired to the real MAC scheduler's transmit-buffer byte count, a physically real quantity two to three orders of magnitude larger in scale. The policy's own decision-facing
input escapes this (a separate path clips the state-facing backlog
term before it reaches the network), which is why the policy's
decisions stayed sensible live even as this diagnostic margin
exploded. Checking 28 seeds of an already-committed prior live campaign on this rig confirms the calibration constant is not simply ``wrong'', since it reproduces a sensible number under normal conditions and only explodes once a real, independently-documented backlog failure regime is entered. Rather than leave the mismatch diagnosed-but-unpatched, we recalibrated using a Little's-Law byte-to-latency conversion against each slice's own latency budget, a physically meaningful quantity rather than an arbitrary rescaling constant, which brings the catastrophic case ($\approx$10.02MB backlog) to $-809$ to $-1{,}187$ (36--53 real seconds of queueing delay against a 45ms budget) rather than $-1{,}002{,}377.5$. This is three orders of magnitude smaller and directly interpretable, without reopening the rank-prediction question already closed above.

\subsection{Heavier Live Load}
\label{sec:results-m8-load}
\prefixflag{Same pre-M41-fix config as sec:results-m8 above. The complete 6-UE collapse reported here may be the floor-starvation bug's own universal failure signature (which onsets independent of UE count or traffic shape, per M41's own root-cause finding) rather than a genuine out-of-distribution generalisation failure. CONFIRM/REVISE/OVERTURN pending M49b-3.}
The anchor above uses 3 UEs and 2 episodes, matching the reproduction budget of the original live sanity check. To ask whether its finding holds under heavier load with real statistical weight rather than a two-episode snapshot, we ran the same seed-900 single-agent-DQN checkpoint for 20 live episodes at 3 UEs and, separately, 20 live episodes at 6 UEs (two per slice), real per-slice UDP traffic throughout, same episode length, 95\% bootstrap CIs (10{,}000 resamples) over the real per-episode values. Fig.~\ref{fig:live-3v6} reports both campaigns directly from their raw logs.

The two loads produced qualitatively different live behaviour from the same frozen checkpoint, not just a quantitative shift. At 3 UEs, block precision was perfect in every one of the 20 episodes (1.000, mmTC-only). At 6 UEs, the same checkpoint blocked nothing at all in any of the 20 episodes; precision is undefined for the entire campaign, not merely low. This is a complete, sustained collapse to always-accept, not an occasional lapse, and it holds throughout a full 100-minute live session rather than clearing after one or two episodes. Mean reward per step is significantly worse at 6 UEs (mean $-8.363$, 95\% CI [$-8.757$, $-7.975$], against 3 UEs' $-7.091$, 95\% CI [$-7.484$, $-6.703$]; the intervals do not overlap). SLA margins worsen substantially and consistently across all three slices (URLLC $+100\%$, eMBB $+100\%$, mmTC $+92\%$ more negative at 6 UEs, each with a bootstrap CI that excludes zero). Doubling live load did not just cost margin on the slices the 3-UE checkpoint left unprotected; it removed the differentiated-shedding behaviour itself for the full duration of this campaign.

\begin{figure}[t]
\centering
\includegraphics[width=0.92\textwidth]{figures/fig7_live_3v6.pdf}
\caption{Live single-gNB comparison, same seed-900 checkpoint, 3 UEs vs.\ 6 UEs, 20 episodes each, real per-slice UDP traffic throughout. (a) SLA margin percent change per slice from 3 to 6 UEs with 95\% bootstrap CI (red: worse) -- all three slices worsen substantially. (b) Mean reward per step, mean $\pm$ 95\% CI with individual episodes shown behind it -- the 6-UE campaign is significantly worse and the intervals do not overlap. Block precision (not plotted, since one side is a single number and the other is undefined) was 1.000 in all 20 episodes at 3 UEs and undefined (zero blocks) in all 20 episodes at 6 UEs.}
\label{fig:live-3v6}
\end{figure}

This comparison remains single-gNB throughout, matching every other live result in this paper: it demonstrates that this checkpoint's differentiated-shedding behaviour, reliable and precise at 3 UEs across 20 independent episodes, does not merely degrade but disappears entirely under sustained heavier load on the hardware this rig actually supports. This is not a multi-gNB coordination claim, and we do not extrapolate it to loads or architectures beyond what was directly measured here.

\subsection{Recalibrating the Simulator to Match Live Congestion}
\label{sec:results-m8-fix}
\prefixflag{This section's headline recalibration result, including its own LIVE confirmation at 6 UEs, also predates the M41 fix. The ``served PRB completely flat against ratio'' reading this section's premise rests on is very likely the \texttt{avg\_prbs\_dl} observability-floor artifact M44-A later documented (\texttt{min\_rbSize=5} is the finest granularity the metric can ever show, not evidence that the ceiling has no real effect), not a genuine physical property of the control surface. Whether the recalibration's own benefit survives once re-run against a fixed floor and honest, non-artifact anchors is exactly what M49b-1 through M49b-4 test. Do not cite this section's numbers as established until that completes.}
The complete 6-UE collapse above raises an obvious question: is it fixable, or an inherent limit of what this checkpoint's architecture can learn? We tested two increasingly targeted offline retraining attempts before finding the actual mechanism, and report both, since neither worked and the reason why is itself informative.

The first attempt doubled the offline simulator's synthetic arrival rate, retrained single-agent DQN from scratch, and ran the result live at 6 UEs for 20 episodes: the checkpoint blocked nothing in any episode, identical to the original. The second attempt doubled the simulator's demand parameter calibrated directly from real single-UE probe measurements rather than an arbitrary synthetic rate. An offline control test first checked whether this lever could even distinguish the original (collapsing) checkpoint's behaviour under 1$\times$ versus 2$\times$ demand: it could not, producing byte-identical block counts and decisions in both conditions. Retraining under it regardless reproduced the same complete live collapse.

Direct instrumentation explained why. We captured the exact 13-dimension input the policy receives at every admission decision -- not the downstream reward, the raw state itself -- both offline and live, via a non-invasive wrapper around the policy's own decision call that adds no new logging path to any frozen file. Offline training only ever exposes the policy to a congestion-level feature in the range 0.03--0.09, regardless of which demand lever is turned, because served capacity in the frozen simulator is capped by the admission ceiling long before demand can push it higher. Live, the same feature sits at 0.23--0.26 at 3 UEs and 0.60--0.69 at 6 UEs -- both already far outside anything the checkpoint was ever trained on, with the checkpoint happening to still discriminate correctly in the range 3 UEs produces and collapsing to always-accept in the range 6 UEs produces. This is an out-of-distribution extrapolation artifact in both directions, not a genuine 3-vs-6-UE competence gap.

A live sweep of the admission ceiling itself, spanning the entire range a trained policy can issue (ratio 1 through 4, the operative bounds in the live config), found served PRB completely flat against ratio for all three slices (urllc/mmtc at 5.00, eMBB at 13.00, $n=20$ polls per point, real traffic throughout) while backlog climbed underneath regardless of the setting. On this hardware, in this regime, the admission-ceiling lever a policy adjusts every step has no measured effect on throughput at all; only the accept/reject decision does, through the backlog relief already built into the simulator's reject-handling path. This value cross-confirms the independently captured state-vector reading above (same order of magnitude, measured two different ways), and reframes recalibration: rather than search for the right ratio-to-PRB curve, there is no such curve to find, and the simulator instead needs served capacity treated as a measured, ratio-independent constant.

We built a new, non-frozen offline environment implementing exactly that: served PRB per slice draws from the real 3-UE and 6-UE anchors above, interpolated by a slowly mean-reverting random walk so a single training run sees the entire realistic congestion range rather than one fixed point, with offered demand scaled by the same interpolation (oversubscribed 1.3$\times$ over served, so backlog genuinely accumulates at both ends of the range, matching live's own behaviour of saturating within minutes regardless of load level). A first version of this environment left offered demand fixed while only served capacity scaled with load -- a real implementation bug that made high-load steps easier, not harder, and still collapsed under it. Correcting this fixed the offline result: across four seeds, three (900, 902, 903; precision 0.99--1.00, zero collapsed episodes out of fifty each) now learn a genuine, consistent differentiated-shedding policy spanning the full congestion range. The fourth (901) still collapses, but this is the same seed that collapsed under the original training environment too, an already-documented, seed-intrinsic single-agent-DQN instability orthogonal to the live-congestion mismatch this recalibration targets.

Taking the recalibrated seed-900 checkpoint live at 6 UEs, 20 episodes, real traffic throughout, reverses the collapse completely: every episode blocked, every block correctly targeting mmTC (precision 1.000, zero collapsed episodes), essentially reproducing the original checkpoint's own successful 3-UE behaviour under the load that previously erased it entirely (Fig.~\ref{fig:live-fix}). Mean reward per step moves back toward the healthy range (mean $-7.816$, 95\% CI [$-8.195$, $-7.442$], against the original checkpoint's $-8.363$ [$-8.757$, $-7.975$] at the same load and $-7.091$ [$-7.484$, $-6.703$] at 3 UEs); this interval only marginally separates from the collapsed checkpoint's own, consistent with this paper's repeated finding that reward is a weaker discriminator on this rig than block precision, which here is unambiguous.

\begin{figure}[t]
\centering
\includegraphics[width=0.92\textwidth]{figures/fig8_live_recalibrated_fix.pdf}
\caption{Three live campaigns, 20 episodes each, real per-slice UDP traffic throughout: the original checkpoint at 3 UEs, the same checkpoint at 6 UEs, and a checkpoint retrained against the recalibrated simulator at 6 UEs. (a) Blocks per episode, mean and range across all 20 episodes -- the original checkpoint's complete silence at 6 UEs sits between two conditions of comparable, healthy block volume. (b) Mean reward per step, mean $\pm$ 95\% CI with individual episodes behind it. Block precision (not plotted; one condition is undefined) was 1.000 for both the original-3UE and recalibrated-6UE conditions and undefined (zero blocks) for the original-6UE condition.}
\label{fig:live-fix}
\end{figure}

This does not establish that recalibrating against a handful of measured anchors generalises beyond the two load levels measured here, or that it would transfer to a different checkpoint architecture; it establishes that this specific collapse was an artifact of a training environment whose congestion never resembled reality, not an inherent ceiling on what this checkpoint could learn, and that correcting the mismatch directly -- rather than retraining under more of the same kind of wrong demand -- produces a policy that stays correctness-aware under the load that previously broke it.

\subsection{Extending the Recalibration to the Topology-Scaling Axis}
\label{sec:results-m27}
The recalibration above targets a single-gNB live collapse; \S\ref{sec:results-m6}'s own topology-scaling campaign was run entirely offline, entirely against the original, uncorrected simulator. Since that simulator's congestion range was shown above to be structurally capped well below reality, we reran the identical campaign (12 seeds, three topologies, three arms, same training/eval budget) against the recalibrated simulator at both scales \S\ref{sec:results-m6} tested, $N{=}19$ and $N{=}7$, to check whether its findings were an artifact of the same miscalibration or hold up once each simulated gNB's congestion actually resembles this rig's measured live range.

At $N{=}19$, they hold up for both of the paper's headline comparison points. Single-agent DQN's collapse rate is statistically unchanged (91.7\% [75.0\%, 100\%] against the original 100\%, $15/15$), and GAT-CTDE's is likewise unchanged (33.3\% [11.1\%, 58.3\%] against the original pooled estimate of 35.4\% [25.5\%, 45.8\%]), with the same qualitative pattern in the non-collapsed seeds -- mostly near-ceiling precision, with a low- or zero-precision seed appearing at every topology. Independent DQN is the one place this recalibration surfaces something new at this scale: the original claim of never fully collapsing ($0/36$) does not survive unchanged, with $3/36$ recalibrated cells showing a genuine total collapse (8.3\% [0.0\%, 25.0\%]) -- a small, real difference we report rather than smooth over, though the recalibrated CI does not exclude the original 0\%.

At $N{=}7$, the picture is more informative than a simple replication. \S\ref{sec:results-m6}'s own $N{=}7$ number was a 3-seed pilot, never resampled at scale because that work judged $N{=}19$'s collapse question the higher-value target -- GAT-CTDE and single-agent DQN both held 1.000 precision in that pilot ($0/3$ collapsed each). A properly powered 12-seed resample tells a different story for one of them: GAT-CTDE's $N{=}7$ collapse rate (33.3\% [8.3\%, 58.3\%]) is now statistically indistinguishable from its own $N{=}19$ rate, suggesting its collapse-proneness is a property of the architecture under real congestion, present from $N{=}7$ onward, not something that only emerges at large $N$. Single-agent DQN is the more striking case: half of the 36 recalibrated cells collapse completely (50.0\% [25.0\%, 75.0\%]), directly contradicting the original pilot's "holds 1.000 at $N{=}7$." We cannot cleanly separate how much of this is the recalibration itself versus the original $n{=}3$ pilot simply never being large enough to see a $\sim$50\% rate -- both are almost certainly true simultaneously, and disentangling them would need an equally sized $N{=}7$ resample against the original simulator, which we did not run (it would not have answered this section's actual question). Independent DQN replicates cleanly at $N{=}7$ (still $0/36$).

\begin{figure}[t]
\centering
\includegraphics[width=\columnwidth]{figures/fig9_m27_scaling_reframe.pdf}
\caption{Collapse rate, original topology-scaling campaign (grey, ClosedLoopKpmSource, \S\ref{sec:results-m6}'s own published values) against this recalibrated resample (teal, RealisticServedKpmSource, 12 seeds $\times$ 3 topologies, bootstrapped over seeds since collapse status correlates within a seed across its topologies). (a) $N{=}7$: the original bars are a single $n{=}3$ pilot point estimate with no CI ever reported, not a zero-width interval. (b) $N{=}19$: both sides are properly powered. Single-agent DQN and GAT-CTDE are statistically unchanged at $N{=}19$; at $N{=}7$, GAT-CTDE now matches its own $N{=}19$ rate and single-agent DQN collapses far more than the original pilot suggested.}
\label{fig:m27-reframe}
\end{figure}

The two scales tell different parts of the same honest story. At $N{=}19$ -- the scale carrying this paper's actual comparative claims -- the recalibration validates rather than retracts \S\ref{sec:results-m6}: GAT-CTDE's advantage over single-agent DQN on collapse resistance was not an artifact of the original simulator's miscalibrated congestion range. At $N{=}7$, the recalibration (or simply the thirteen-fold larger sample that came with it) reveals real complexity a 3-seed pilot could never have shown: single-agent DQN is not reliably safe at small scale either, and GAT-CTDE's collapse-proneness does not wait for a large cluster to appear. Two concrete corrections follow: independent DQN's "never fully collapses" claim softens to "rarely, not never" (consistent at both scales, $0/36$ and $3/36$), and any reading of $N{=}7$ as a collapse-free scale for single-agent DQN should be retracted -- it was never well-supported, only under-sampled.

This section's own live anchor uses the SLA-only reward and, at 20--24 episodes per campaign, is powered to detect a collapse but not to carry a formal significance claim. \prefixinline{that live anchor is pre-M41-fix, re-collection pending (docs/PAPER5\_M48b\_provenance\_audit.csv)} A companion submission~\cite{cacs26-paper} separately measures the E2 control loop itself as negligible overhead (median round-trip 0.57\,ms) from a policy under 20,000 parameters, on the same rig.

\section{Conclusion}
\label{sec:conclusion}
We set out to test whether a standard evaluation metric can be
trusted in multi-agent, privacy-constrained, disruption-prone slice
admission control, and found it cannot go unchecked. A standard SLA-compliance metric, an unguarded append-only log, and an unseeded weight-initialisation path together hid three distinct failures, namely a training collapse the compliance metric scored no worse than correct behaviour, a hidden eval-data contamination bug, and a same-seed-only reproducibility gap, none of which a standard pipeline without this paper's correctness-aware metrics and reproduction-replication discipline would have caught.

GAT-CTDE, the multi-agent architecture that served as this evaluation's substrate, delivers a decisive, statistically significant improvement over an independent-learner ablation on the correctness-aware reward metric, though its edge over our existing single-agent baseline specifically is small and not decisive at this sample size. This is an honest, non-inflated architectural finding, not the headline result this paper argues evaluation practice should not be built around alone.

Extending to federated training with differential privacy, federation
adds no measurable reward cost, and privacy is costly only past a real
threshold in decision precision, not a smooth trade-off; FedProx earns
no measurable dividend under genuine per-gNB load heterogeneity at any
tested strength, traced to real per-round client drift being too small
relative to the loss scale to matter, even at ten times the round
length. Under mid-episode gNB dropout, demand spikes, and agent churn, dropout and churn produce a real, accelerating cost with severity, while demand spikes do not threshold at all on the metric that matters. Cross-checking every result against an independent seed sample, not just the same seeds reproduced again, caught one architectural claim (that independent DQN is immune to agent churn) that a single sample had made look confirmed but was not, and it is retracted here rather than reported.

Scaling the cluster to 19 gNBs, single-agent DQN's collapse becomes
total and GAT-CTDE's remains only partial, a real cluster-size-specific
effect; GAT-CTDE's own collapse rate, read at two disagreeing values
across two small samples (31\% vs.\ 78\%), was resolved rather than
left open with a third, independently-drawn sample, converging on
35.4\% [25.5\%, 45.8\%] across 64 seeds. Finally, the offline-trained
admission-control decision logic transferred to a real single-gNB
testbed essentially unchanged, while its SLA-margin metric diverged by
orders of magnitude for reasons traced to a concrete unit mismatch
rather than left unexplained.

Two limits bound every claim in this paper. First, every multi-gNB, federated, and disruption result is offline-only, since this rig cannot yet sustain a live multi-gNB test, so whether these rankings would survive contact with a live, multi-gNB testbed remains untested, not just under-tested. A companion submission~\cite{cacs26-paper} establishes that this rig's live E2 control loop is itself well-behaved for a single gNB, negligible round-trip latency, a lightweight policy, but whether that same loop stays well-behaved once it is coordinating $N>1$ gNBs live is exactly the untested step this limit names, not something either paper yet answers.
Second, this paper's federated scenario is a synthetic privacy setting (one operator's cluster, not genuine cross-organisation data separation), so we lead with the operational benefit of periodic, not per-step, synchronisation for disaggregated RAN Intelligent Controllers rather than a privacy motivation this scale cannot substantiate. Together, these results argue that admission-control
evaluation in multi-agent, privacy-constrained, disruption-prone
settings needs metrics tied to the actual training objective, not only
to a downstream SLA proxy that can reward the wrong policy, and that
architectural claims need more than one seed sample before they are
trusted.

\backmatter

\bmhead{Acknowledgements}
This work was supported by the Fundamental Research Grant Scheme
(FRGS) of the Ministry of Higher Education (MOHE), Malaysia, under
grant FRGS/1/2024/ICT06/TAYLOR/02/1.

\section*{Declarations}

\noindent The following statements should be understood to apply to
this manuscript unless otherwise indicated.

\begin{itemize}
\item \textbf{Funding:} This work was supported by the Fundamental Research Grant Scheme (FRGS), Ministry of Higher Education (MOHE), Malaysia, under grant number FRGS/1/2024/ICT06/TAYLOR/02/1.
\item \textbf{Competing interests:} The authors declare no competing interests.
\item \textbf{Ethics approval and consent to participate:} Not applicable, since this study involves no human participants, animal subjects, or identifiable patient data; all reported results are either offline simulation campaigns or measurements from a laboratory radio testbed.
\item \textbf{Consent for publication:} Not applicable, since this manuscript contains no personal data or identifying information about any individual.
\item \textbf{Data availability:} The experiment logs, checkpoints, and campaign result files underlying this paper's findings are available from the corresponding author on reasonable request.
\item \textbf{Materials availability:} Not applicable, since this study involves no physical or biological materials.
\item \textbf{Code availability:} The framework and experiment scripts used to produce this paper's results are available from the corresponding author on reasonable request.
\item \textbf{Author contribution:} M.\ P.\ Kunasegran: conceptualisation, methodology, software, investigation, writing (original draft). W.\ L.\ Pang: supervision, writing (review and editing). S.\ K.\ Phang: supervision, writing (review and editing).
\end{itemize}

\bibliography{refs}

\end{document}
